โครงงานนี้พัฒนาระบบเว็บสำหรับลดความเสี่ยงจากการเปิดเผยตัวตนในภาพและวิดีโอ โดยใช้ปัญญาประดิษฐ์ตรวจจับใบหน้า จัดกลุ่มใบหน้าของบุคคลเดียวกัน และปกปิดพื้นที่ที่เลือกด้วยเอฟเฟกต์ Gaussian Blur, Pixelation หรือ Black Box ผู้ใช้สามารถอัปโหลดไฟล์ภาพและวิดีโอ เลือกบุคคลที่จะปกปิดจาก Face Map ตรวจผลลัพธ์ก่อนส่งออก และกำหนดพื้นที่ปกปิดด้วยตนเองสำหรับวัตถุหรือข้อมูลอื่นที่ระบบตรวจจับไม่ได้

ในการพัฒนารอบปัจจุบัน ส่วนหน้าพัฒนาด้วย Next.js, React และ TypeScript ส่วนหลังบ้านพัฒนาด้วย FastAPI, SQLAlchemy และ SQLite โดยใช้ YOLO สำหรับตรวจจับใบหน้า InsightFace/ArcFace สำหรับสร้างเวกเตอร์ลักษณะใบหน้า OpenCV สำหรับประมวลผลเฟรม และอัลกอริทึมจับคู่กรอบแบบ IoU ร่วมกับ Hungarian Algorithm สำหรับติดตามตำแหน่งใบหน้าระหว่างเฟรม ระบบบันทึกข้อมูลวิดีโอ กรอบใบหน้า บุคคลที่ตรวจพบ การเลือกเบลอ และกรอบปกปิดที่ผู้ใช้วาดไว้เพื่อใช้สร้างตัวอย่างและไฟล์ผลลัพธ์

ผลการตรวจสอบโครงสร้างโปรแกรมพบว่าฟังก์ชันหลักตั้งแต่การอัปโหลด การตรวจจับและจัดกลุ่มใบหน้า การเลือกปกปิด การวาดกรอบด้วยตนเอง การติดตามกรอบ การแสดงตัวอย่าง และการส่งออกไฟล์ได้รับการพัฒนาแล้ว อย่างไรก็ตาม การประเมินความแม่นยำและเวลาในชุดข้อมูลมาตรฐาน รวมถึงการทดสอบแบบ end-to-end ในสภาพแวดล้อมที่ตั้งค่าแบบจำลองและ Cloudinary ครบถ้วน ยังเป็นงานที่ต้องดำเนินการต่อไป
